that are abelian, as follows: for @@M0@@, we associate to @@M1@@ the set of formal series @@M2@@ satisfying the following conditions: if @@M3@@, then
\sgaStableTarget{sga6.e05.eq.1-1-1}@@M4@@
This set will be denoted by @@M5@@, or by @@M6@@ when no confusion is possible; it is clear that it is an abelian group under multiplication of formal series, and that this does define a functor
@@M7@@
For @@M8@@ endowed with the coarse filtration, @@M9@@ is the group of formal series of augmentation @@M10@@, which will also be denoted by @@M11@@.

\sgaStableTarget{sga6.e05.lem.1-2}\textbf{Lemma 1.2.} \uline{The functor} @@M12@@ \uline{is representable}.

Let @@M13@@ be a family of indeterminates; we assign the weight @@M14@@ to @@M15@@, and filter the ring @@M16@@ by weight with respect to the @@M17@@. We then define a bijection:
@@M18@@
by associating to every series @@M19@@ the homomorphism @@M20@@ determined by @@M21@@. This homomorphism is compatible with the filtrations, and we thus obtain an isomorphism of functors.

Put @@M22@@; for every filtered @@M23@@-algebra @@M24@@ and every series @@M25@@, there exists exactly one homomorphism
